\documentclass[11pt]{article}
\usepackage{cs152}

%%%%%%%%%%%%%%%%%%%% name/id
\rfoot{\small Brian Park | 3033045069}


%%%%%%%%%%%%%%%%%%%% Course/HW info
\newcommand*{\instr}{John Wawrzynek}
\newcommand*{\term}{Spring 2022}
\newcommand*{\coursenum}{CS 152}
\newcommand*{\coursename}{Computer Architecture and Engineering}
\newcommand*{\hwnum}{1}

\rhead{\LARGE   \fontfamily{lmdh}\selectfont	HW \hwnum}

\lfoot{\small \coursenum, \term, HW \hwnum}

%%%%%%%%%%%%%%%%%%%%%%%%%%%%%% Document Start %%%%%%%%%%%%%%%%%
\begin{document}

%%%%%%%%%%%%%%%%%%%%%%%%%%%%%%%%%%%%%%%%%%%%%%%%%%%%%%%%%%%%%%%%%%%%%%%%%%%%%%%%%%%%%%%%
% Question 3.4
%%%%%%%%%%%%%%%%%%%%%%%%%%%%%%%%%%%%%%%%%%%%%%%%%%%%%%%%%%%%%%%%%%%%%%%%%%%%%%%%%%%%%%%%
\Question{CISC, RISC, Accumulator, and Stack: Comparing ISAs}

\lstset{language=[RISC-V]{Assembler}}
\begin{lstlisting}
loop:  test %ecx,%ecx 
       jz done
       mov %eax,%ebx 
       shr $31,%ebx 
       shl $1,%eax 
       or %ebx,%eax d
       ec %ecx
       jmp loop 
done:  ...
\end{lstlisting}

\begin{Parts}
	\Part In total there are 15 bytes in the program. There needs to be 79 bytes of instructions needed to be fetched. We have to be aware that the loop condition is fetched $n + 1$ iterations, so the total number bytes fetched is $(6 * 4) + (5 * 11) = 79$ by accounting the number of iterations in the loop condition and loop body. In the program, there is no memory stored or loaded, at least in the program given. Thus 0 bytes of data loaded or stored.
	\Part Here lies the translation from x86 to RISC-V assembly:
\lstset{language=[RISC-V]{Assembler}}
\begin{lstlisting}
loop:  beq a0, x0, done
       mov t0, a1
       srli t0, t0, 31
       slli a1, a1, 1
       or a1, t0, a1
       addi a0, a0, -1
       j loop
done:  ...
\end{lstlisting}
	The instruction set between 32 bit and 64 bit ISA should remain the same bit width, so all we would need to do is multiply how many instructions are in the converted assembly code by 4 bytes. Thus, there is 1 instruction for the branch condition and 6 instruction for the loop body. $(6 * 4) + (5 * 24) = 144$ bytes of instructions fetched. Again, there is no loading or storing, so 0 bytes of memory is accessed/transferred. 
	
	\Part Here lies the translation from x86 to stack assembly:
	
\lstset{language=[RISC-V]{Assembler}}
\begin{lstlisting}
       pop 0x8000
loop:  push 0x8000
       beqz done
       pop 0x8004
       push 0x8004
       sign
       push 0x8004
       shl
       or
       push 0x8000
       dec
	   jump loop
done:  ...
\end{lstlisting}

Push, pop, and branch instructions are the most expensive and take 3 bytes. Other instructions are 1 byte. The program itself has $(7 * 3) + (4 * 1) = 25$ bytes. The number of bytes that needs to be fetched is 119 bytes. The number of bytes loaded depends on the number of pushes, so there are $(4 * 6) + (4 * 3 * 5) = 84$ bytes loaded from memory. The number of bytes stored depends on the number of pops, so there are  $4 + (4 * 5) = 24$ bytes stored to memory.

\end{Parts}


\newpage
\Question{CISC vs RISC}
\begin{Parts}
	\Part Assume that compiler technology is poor, and therefore your users are far more apt to write all of their code in assembly. A \textbf{RISC} ISA would be best appreciated by these programmers. \\ \\
	Although CISC provides a more complex ISA, RISC may be more suitable towards the programmer. For example, RISC-V has 47 instructions. It may take more instructions to write a program compared to CISC, but at least it's better than trying to search what the correct instruction to use is in the thousands of instructions that CISC, such as x86 provides.
	\Part You desire to make compilers better at targeting your \textit{yet-to-be-designed} machine. Therefore, you choose a \textbf{RISC} ISA, as it would be easiest for a compiler to target, thus allowing your users to write code in higher-level languages like C and Fortran and raise their productivity. \\ \\
	Because a CISC ISA may have a lot of microcoded subroutines and microinstructions, there maybe complicated conversion to a co
	\Part Assume that CPU logic is fast, very fast, while instruction fetch accesses are at least 10x slower (suppose you are the lead architect of the “709”). Which ISA style do you choose as a best match for the hardware’s limitations? \\ \\
	\textbf{CISC}
	\Part Starting with a clean slate in the year 2021 (area/logic/memory is cheap), you think that a \textbf{RISC} ISA that would lend itself best to a very high performance processor (e.g., high frequency, highly pipelined).

	
\end{Parts}




\end{document};